% !TeX program = lualatex
\documentclass[12pt,aspectratio=169]{beamer}
\usepackage[utf8]{inputenc}
\usepackage[T1]{fontenc}
\usepackage{amsmath,amsfonts,amssymb}
\usepackage{graphicx}
\usepackage{xcolor}
\usepackage{hyperref}
\usepackage{anyfontsize}
\usepackage{lmodern}
\usepackage{longtable}
\usepackage{array}
\usepackage{booktabs}

% ====== EXTREME FONT REDUCTION ======
\renewcommand{\tiny}{\fontsize{3.5}{4.5}\selectfont}
\renewcommand{\scriptsize}{\fontsize{4.5}{5.5}\selectfont}
\renewcommand{\footnotesize}{\fontsize{5.5}{6.5}\selectfont}
\renewcommand{\small}{\fontsize{6.5}{7.5}\selectfont}
\renewcommand{\normalsize}{\fontsize{7.5}{8.5}\selectfont}
\renewcommand{\large}{\fontsize{8.5}{9.5}\selectfont}
\renewcommand{\Large}{\fontsize{9.5}{10.5}\selectfont}
\renewcommand{\LARGE}{\fontsize{10.5}{12}\selectfont}
\renewcommand{\huge}{\fontsize{11.5}{13.5}\selectfont}
\renewcommand{\Huge}{\fontsize{13}{16}\selectfont}

% ====== COLORS ======
\definecolor{redcorrect}{RGB}{200,30,30}
\definecolor{bluecorrect}{RGB}{30,80,200}
\definecolor{greencheck}{RGB}{0,150,50}
\definecolor{lightgray}{RGB}{245,245,245}
\definecolor{darkgray}{RGB}{40,40,40}
\definecolor{softyellow}{RGB}{255,248,200}
\definecolor{deepblue}{RGB}{20,60,150}
\definecolor{darkred}{RGB}{150,20,20}
\definecolor{orangewarning}{RGB}{255,140,0}
\definecolor{correctgreen}{RGB}{0,120,40}
\definecolor{trapcolor}{RGB}{180,80,0}
\definecolor{purpleconcept}{RGB}{120,50,180}
\definecolor{tealhard}{RGB}{0,120,120}

\setbeamercolor{title}{fg=white,bg=redcorrect}
\setbeamercolor{frametitle}{fg=white,bg=redcorrect}
\setbeamercolor{block title}{fg=white,bg=bluecorrect}
\setbeamercolor{block body}{fg=darkgray,bg=lightgray}
\setbeamercolor{normal text}{fg=darkgray}
\usecolortheme{default}

% ====== CUSTOM COMMANDS ======
\newcommand{\correct}[1]{\textcolor{greencheck}{\textbf{\checkmark}} #1}
\newcommand{\keypoint}[1]{\textcolor{deepblue}{\textbf{#1}}}
\newcommand{\hardconcept}[1]{\textcolor{darkred}{\textbf{#1}}}
\newcommand{\examtraps}[1]{\textcolor{trapcolor}{\textbf{⚠ TRAP:}} #1}
\newcommand{\recallbox}[1]{\fcolorbox{deepblue}{blue!5}{\parbox{0.88\textwidth}{\centering \textbf{REMEMBER:} #1}}}
\newcommand{\stepbox}[1]{%
	\begin{center}
		\fcolorbox{darkgray}{white}{%
			\parbox{0.92\textwidth}{\vspace{0.03cm} #1 \vspace{0.03cm}}%
		}
	\end{center}
}
\newcommand{\answerbox}[1]{\fcolorbox{correctgreen}{green!8}{\parbox{0.88\textwidth}{\centering \textcolor{correctgreen}{\textbf{\checkmark\ CORRECT:}} #1}}}
\newcommand{\wronganswer}[1]{\textcolor{redcorrect}{\textbf{✘}} #1}
\newcommand{\trapbox}[1]{\fcolorbox{trapcolor}{orange!8}{\parbox{0.88\textwidth}{\centering \textcolor{trapcolor}{\textbf{⚠ EXAM TRAP:}} #1}}}
\newcommand{\keyrule}[1]{\textcolor{tealhard}{\textbf{🔑}} #1}

\begin{document}
	
	% ================================================================
	% SLIDE 1: TITLE
	% ================================================================
	\begin{frame}
		\centering
		\vspace{0.5cm}
		{\fontsize{18}{22}\selectfont \textbf{ALL CAMS DOMAINS}}\\[0.15cm]
		{\fontsize{14}{17}\selectfont \textbf{HARD QUESTIONS \& EXAM TRAPS FROM ALL IMAGES}}\\[0.2cm]
		{\fontsize{9}{11}\selectfont \textit{Everything You Get Wrong — Explained in Detail}}\\[0.1cm]
		\rule{4cm}{0.5pt}\\[0.1cm]
		{\fontsize{7}{9}\selectfont Domain A: VASP · Crypto · FATF Recs · Sanctions · DORA · Investment Banking}\\[0.05cm]
		{\fontsize{7}{9}\selectfont Domain B: MLAT · PPP · FATF IOs · AMLA · Three Lines · EWRA}\\[0.05cm]
		{\fontsize{7}{9}\selectfont Domain C: AML Policy · Board Oversight · PEP EDD · SOW/SOF · SAR/STR}\\[0.05cm]
		{\fontsize{7}{9}\selectfont Domain D: Transaction Monitoring · OSINT · Data Lineage · AI/ML · Clustering}
		\vspace{0.3cm}
	\end{frame}
	
	% ================================================================
	% SLIDE 2: VASP - FATF REC 15 - THE CRYPTO TRAP
	% ================================================================
	\begin{frame}
		\frametitle{\fontsize{11}{13}\selectfont \textbf{VASP — FATF Rec 15 \& MiCA — The Crypto Trap}}
		\begin{block}{\fontsize{7}{9}\selectfont \textbf{The Scenario — CAMS Exam Favorite}}
			\scriptsize
			A startup in Germany plans to issue a new crypto token and market it throughout the EU.
		\end{block}
		\vspace{0.02cm}
		\begin{block}{\fontsize{7}{9}\selectfont \textbf{The Hard Question}}
			\scriptsize
			What step must it take under MiCA regulations?
		\end{block}
		\vspace{0.02cm}
		\answerbox{\large \textbf{Publish a crypto-asset white paper and notify the national authority.}}
		\vspace{0.02cm}
		\begin{block}{\fontsize{7}{9}\selectfont \textbf{Why This Is a TRAP}}
			\scriptsize
			\begin{itemize}
				\item \wronganswer{Seek regulatory approval in each EU country} — \textbf{NO.} MiCA is \textbf{harmonized} — one notification, not each country.
				\item \wronganswer{Register with FATF} — \textbf{NO.} FATF is standard-setter, not regulator.
				\item \textbf{MiCA (Markets in Crypto-Assets)} = EU crypto regulation.
				\item \textbf{White paper requirement:} Must publish detailed disclosure document.
				\item \textbf{Notification:} Notify the national competent authority in home country.
				\item \textbf{Passporting:} Once approved, can operate across all EU countries.
				\item \textbf{FATF Rec 15:} VASPs must be licensed/registered and apply AML/CFT.
			\end{itemize}
		\end{block}
		\vspace{0.02cm}
		\recallbox{\scriptsize \textbf{Key Rule:} MiCA = \textbf{white paper + one national notification}. Passporting across EU.}
	\end{frame}
	
	% ================================================================
	% SLIDE 3: AUTHENTICATION TECHNOLOGIES - THE BIOMETRICS TRAP
	% ================================================================
	\begin{frame}
		\frametitle{\fontsize{11}{13}\selectfont \textbf{Authentication Technologies — The Biometrics Trap}}
		\begin{block}{\fontsize{7}{9}\selectfont \textbf{The Hard Question — CAMS Exam Favorite}}
			\scriptsize
			Which technologies enhance authentication and security in financial systems? (Select Three.)
		\end{block}
		\vspace{0.02cm}
		\answerbox{\large \textbf{1. Multi-factor authentication. 2. Physiological biometrics. 3. Behavioral biometrics.}}
		\vspace{0.02cm}
		\begin{block}{\fontsize{7}{9}\selectfont \textbf{Why This Is a TRAP}}
			\scriptsize
			\begin{itemize}
				\item \wronganswer{Password biometrics} — \textbf{NOT} a recognized category. Passwords are \textbf{not} biometrics.
				\item \textbf{Multi-factor authentication (MFA)} = multiple factors (something you know, have, are).
				\item \textbf{Physiological biometrics} = physical traits (fingerprint, facial recognition, iris scan, voice).
				\item \textbf{Behavioral biometrics} = behavioral patterns (typing rhythm, mouse movements, gait).
				\item \textbf{Password biometrics} = \textbf{trick answer} — there is no such thing.
				\item \textbf{Deeepfake detection:} Tools that detect light/skin texture inconsistencies and liveness prompts (blink, smile).
			\end{itemize}
		\end{block}
		\vspace{0.02cm}
		\recallbox{\scriptsize \textbf{Key Rule:} MFA + physiological biometrics + behavioral biometrics = correct. "Password biometrics" = trap.}
	\end{frame}
	
	% ================================================================
	% SLIDE 4: MLAT - THE ADMISSIBILITY TRAP (DOMAIN B)
	% ================================================================
	\begin{frame}
		\frametitle{\fontsize{11}{13}\selectfont \textbf{MLAT — The Admissibility Trap}}
		\begin{block}{\fontsize{7}{9}\selectfont \textbf{The Scenario — CAMS Exam Favorite}}
			\scriptsize
			Two countries, A and B, are establishing an MLAT. Country A has \textbf{data protection rules} preventing sharing of client identities. Country B bases investigations on \textbf{beneficial ownership disclosures}.
		\end{block}
		\vspace{0.02cm}
		\begin{block}{\fontsize{7}{9}\selectfont \textbf{The Hard Question}}
			\scriptsize
			What challenge is most likely to arise in implementing the MLAT?
		\end{block}
		\vspace{0.02cm}
		\answerbox{\large \textbf{Evidence exchanged under the MLAT will NOT be admissible in Country B's legal system.}}
		\vspace{0.02cm}
		\begin{block}{\fontsize{7}{9}\selectfont \textbf{Why This Is a TRAP}}
			\scriptsize
			\begin{itemize}
				\item \wronganswer{Country A's data protection rules will prevent sharing} — This is \textbf{obvious}. The \textbf{real} issue is \textbf{admissibility}, not sharing.
				\item \textbf{Admissibility} is the \textbf{legal standard} for evidence in court.
				\item Even if Country A \textbf{shares} the data, Country B's courts may \textbf{reject} it if it doesn't meet their evidentiary standards.
				\item MLATs must address both \textbf{sharing} AND \textbf{admissibility}.
				\item This is why \textbf{chain of custody} and \textbf{evidence handling} are critical in MLATs.
			\end{itemize}
		\end{block}
		\vspace{0.02cm}
		\recallbox{\scriptsize \textbf{Key Rule:} MLAT challenge is \textbf{admissibility}, not just sharing. Evidence must be legally usable.}
	\end{frame}
	
	% ================================================================
	% SLIDE 5: PPP - THE REAL-TIME INTELLIGENCE TRAP (DOMAIN B)
	% ================================================================
	\begin{frame}
		\frametitle{\fontsize{11}{13}\selectfont \textbf{PPP — The Real-Time Intelligence Trap}}
		\begin{block}{\fontsize{7}{9}\selectfont \textbf{The Scenario — CAMS Exam Trap}}
			\scriptsize
			An FIU wants to use the PPP model for \textbf{real-time intelligence sharing} to disrupt a terrorist financing network.
		\end{block}
		\vspace{0.02cm}
		\begin{block}{\fontsize{7}{9}\selectfont \textbf{The Hard Question}}
			\scriptsize
			Which approach is most aligned with PPP goals?
		\end{block}
		\vspace{0.02cm}
		\answerbox{\large \textbf{Establishing a formal platform with direct two-way communication on specific threats.}}
		\vspace{0.02cm}
		\begin{block}{\fontsize{7}{9}\selectfont \textbf{Why This Is a TRAP}}
			\scriptsize
			\begin{itemize}
				\item \wronganswer{Anonymized, aggregated format after threat is resolved} — This is \textbf{after-the-fact reporting}, NOT real-time sharing.
				\item \textbf{PPP is about DISRUPTION}, not post-mortem analysis.
				\item \textbf{Direct two-way communication} = real-time = can \textbf{prevent} crimes.
				\item Anonymized reporting \textbf{protects privacy} but \textbf{doesn't enable disruption}.
				\item \textbf{Examples:} UK JMLIT, Australia Fintel, Singapore ACIP.
			\end{itemize}
		\end{block}
		\vspace{0.02cm}
		\recallbox{\scriptsize \textbf{Key Rule:} PPPs are about \textbf{real-time} intelligence sharing to \textbf{disrupt} active threats.}
	\end{frame}
	
	% ================================================================
	% SLIDE 6: EARLY EU DIRECTIVES - THE TRANSPOSITION TRAP (DOMAIN B)
	% ================================================================
	\begin{frame}
		\frametitle{\fontsize{11}{13}\selectfont \textbf{EU Directives — The Transposition Trap}}
		\begin{block}{\fontsize{7}{9}\selectfont \textbf{The Hard Question — CAMS Exam Favorite}}
			\scriptsize
			Which best describes a major shortcoming of early European AML Directives that prompted later revisions?
		\end{block}
		\vspace{0.02cm}
		\answerbox{\large \textbf{Failure to ensure timely and complete transposition by all member states.}}
		\vspace{0.02cm}
		\begin{block}{\fontsize{7}{9}\selectfont \textbf{Why This Is a TRAP}}
			\scriptsize
			\begin{itemize}
				\item \wronganswer{Lack of precise definitions for predicate offenses} — \textbf{NOT} the main shortcoming.
				\item \textbf{Directives require transposition} into national law.
				\item Some countries \textbf{delayed} transposition or did it \textbf{poorly}.
				\item \textbf{Consequence:} Regulatory arbitrage — criminals exploited weak AML in certain EU countries.
				\item \textbf{Lesson:} Harmonization requires \textbf{enforcement} of transposition.
				\item \textbf{3rd Directive (2005):} Incorporated FATF Recs.
				\item \textbf{4th Directive (2015):} Required BO registers.
				\item \textbf{5th Directive (2018):} Extended to crypto/prepaid.
				\item \textbf{6th Directive (2020):} Harmonized offenses and penalties.
			\end{itemize}
		\end{block}
		\vspace{0.02cm}
		\recallbox{\scriptsize \textbf{Key Rule:} Early EU directives failed due to \textbf{poor transposition} and \textbf{regulatory arbitrage}.}
	\end{frame}
	
	% ================================================================
	% SLIDE 7: AMLA - THE SUPERVISORY AUTHORITY TRAP (DOMAIN B)
	% ================================================================
	\begin{frame}
		\frametitle{\fontsize{11}{13}\selectfont \textbf{AMLA — The Supervisory Authority Trap}}
		\begin{block}{\fontsize{7}{9}\selectfont \textbf{The Scenario — CAMS Exam Favorite}}
			\scriptsize
			An EU-based fintech operates online account services in \textbf{5 member states} and has been flagged for non-compliance by \textbf{more than one regulator}.
		\end{block}
		\vspace{0.02cm}
		\begin{block}{\fontsize{7}{9}\selectfont \textbf{The Hard Question}}
			\scriptsize
			Which supervisory authority is most likely to assume oversight?
		\end{block}
		\vspace{0.02cm}
		\answerbox{\large \textbf{The AMLA, due to the firm's cross-border nature and risk profile.}}
		\vspace{0.02cm}
		\begin{block}{\fontsize{7}{9}\selectfont \textbf{Why This Is a TRAP}}
			\scriptsize
			\begin{itemize}
				\item \wronganswer{National regulator of headquarters} — This is \textbf{traditional} supervision, but AMLA changes this.
				\item \wronganswer{National regulator with most customers} — Not a valid basis.
				\item \textbf{AMLA (Anti-Money Laundering Authority)} = new EU body (established 2024).
				\item \textbf{Direct supervision:} AMLA directly supervises selected \textbf{cross-border financial institutions}.
				\item Flagged by \textbf{multiple regulators} = cross-border risk = AMLA jurisdiction.
				\item \textbf{Coordination:} AMLA coordinates national supervisors for consistent supervision.
			\end{itemize}
		\end{block}
		\vspace{0.02cm}
		\recallbox{\scriptsize \textbf{Key Rule:} AMLA = \textbf{centralized supervision} for cross-border, high-risk EU institutions.}
	\end{frame}
	
	% ================================================================
	% SLIDE 8: DORA - THE CENTRALIZED OVERSIGHT TRAP
	% ================================================================
	\begin{frame}
		\frametitle{\fontsize{11}{13}\selectfont \textbf{DORA — The Centralized Oversight Trap}}
		\begin{block}{\fontsize{7}{9}\selectfont \textbf{The Hard Question — CAMS Exam Favorite}}
			\scriptsize
			A cross-border bank operating in several EU countries is reviewing its digital risk governance. What action best aligns with DORA?
		\end{block}
		\vspace{0.02cm}
		\answerbox{\large \textbf{Centralizing oversight of ICT risk management across all member-state branches.}}
		\vspace{0.02cm}
		\begin{block}{\fontsize{7}{9}\selectfont \textbf{Why This Is a TRAP}}
			\scriptsize
			\begin{itemize}
				\item \wronganswer{Relying on each national regulator's cybersecurity requirements} — \textbf{Insufficient}. DORA mandates \textbf{pan-European} oversight.
				\item \textbf{Centralized} = unified governance across all EU operations.
				\item \textbf{Not} just compliance with each country's rules.
				\item Requires \textbf{single} ICT risk governance framework for the entire group.
				\item \textbf{Key requirements:}
				\begin{enumerate}
					\item ICT Risk Governance Framework (unified)
					\item ICT Risk Assessments
					\item Incident Reporting
					\item Third-Party ICT Provider Management
					\item Resilience Testing (threat-led penetration testing)
					\item Information Sharing
				\end{enumerate}
			\end{itemize}
		\end{block}
		\vspace{0.02cm}
		\recallbox{\scriptsize \textbf{Key Rule:} DORA = \textbf{centralized} ICT oversight. National rules are \textbf{not enough}.}
	\end{frame}
	
	% ================================================================
	% SLIDE 9: TCSP - THE CONFIDENTIALITY TRAP
	% ================================================================
	\begin{frame}
		\frametitle{\fontsize{11}{13}\selectfont \textbf{TCSP — The Confidentiality Trap}}
		\begin{block}{\fontsize{7}{9}\selectfont \textbf{The Scenario — CAMS Exam Favorite}}
			\scriptsize
			A TCSP firm is reviewing whether to accept a client who wants to incorporate in a low-tax Caribbean jurisdiction. The client requests confidentiality protections for all shareholders. A connection is identified to industries marked high-risk in the NRA.
		\end{block}
		\vspace{0.02cm}
		\begin{block}{\fontsize{7}{9}\selectfont \textbf{The Hard Question}}
			\scriptsize
			What should the TCSP response be?
		\end{block}
		\vspace{0.02cm}
		\answerbox{\large \textbf{Decline immediately unless full beneficial owner information and risk justifications are disclosed and verified.}}
		\vspace{0.02cm}
		\begin{block}{\fontsize{7}{9}\selectfont \textbf{Why This Is a TRAP}}
			\scriptsize
			\begin{itemize}
				\item \wronganswer{Refer to offshore local counsel and detach from compliance responsibility} — \textbf{NO.} TCSP retains compliance responsibility.
				\item \wronganswer{Accept with standard CDD} — \textbf{NO.} High-risk indicators require EDD or decline.
				\item \textbf{Confidentiality requests + high-risk industries = significant red flags.}
				\item \textbf{FATF Rec 24/25:} BO identification is mandatory.
				\item \textbf{Offloading responsibility} to local counsel = \textbf{not} a valid risk mitigation strategy.
				\item \textbf{Alternative:} Refer to internal high-risk committee for EDD before acceptance.
			\end{itemize}
		\end{block}
		\vspace{0.02cm}
		\recallbox{\scriptsize \textbf{Key Rule:} TCSP cannot outsource AML responsibility. Decline unless BO info is disclosed and verified.}
	\end{frame}
	
	% ================================================================
	% SLIDE 10: PEP EDD REQUIREMENTS - THE DOMESTIC/FOREIGN TRAP
	% ================================================================
	\begin{frame}
		\frametitle{\fontsize{11}{13}\selectfont \textbf{PEP EDD — The Domestic/Foreign Trap}}
		\begin{block}{\fontsize{7}{9}\selectfont \textbf{The Hard Question — CAMS Exam Favorite}}
			\scriptsize
			Which controls would be suitable when onboarding a PEP? (Select Three.)
		\end{block}
		\vspace{0.02cm}
		\answerbox{\large \textbf{1. Conducting source of wealth checks. 2. Running checks against sanctions and watchlists. 3. Mandating senior management-level approval.}}
		\vspace{0.02cm}
		\begin{block}{\fontsize{7}{9}\selectfont \textbf{Why This Is a TRAP}}
			\scriptsize
			\begin{itemize}
				\item \wronganswer{Standard CDD only} — \textbf{NO.} PEPs require EDD.
				\item \textbf{Foreign PEPs:} Mandatory EDD (FATF Rec 12).
				\item \textbf{Domestic PEPs:} Risk-based — apply EDD if high risk.
				\item \textbf{SOW/SOF:} Must verify source of wealth and source of funds.
				\item \textbf{Sanctions/watchlists:} Must screen against all lists.
				\item \textbf{Senior approval:} Senior management must approve the relationship.
				\item \textbf{Ongoing monitoring:} Enhanced transaction monitoring and regular reviews.
			\end{itemize}
		\end{block}
		\vspace{0.02cm}
		\recallbox{\scriptsize \textbf{Key Rule:} Foreign PEPs = \textbf{mandatory EDD}. Domestic PEPs = \textbf{risk-based EDD}.}
	\end{frame}
	
	% ================================================================
	% SLIDE 11: SOW vs SOF - THE CRITICAL DISTINCTION
	% ================================================================
	\begin{frame}
		\frametitle{\fontsize{11}{13}\selectfont \textbf{SOW vs SOF — The Critical Distinction}}
		\begin{block}{\fontsize{7}{9}\selectfont \textbf{The Hard Question — CAMS Exam Favorite}}
			\scriptsize
			What is the difference between Source of Wealth and Source of Funds?
		\end{block}
		\vspace{0.02cm}
		\begin{block}{\fontsize{7}{9}\selectfont \textbf{The Critical Distinction}}
			\scriptsize
			\begin{tabular}{|p{0.35\textwidth}|p{0.61\textwidth}|}
				\hline
				\keypoint{Source of Wealth (SOW)} & \keypoint{Source of Funds (SOF)} \\
				\hline
				\textit{How} did you get your money? & \textit{Where} did this specific money come from? \\
				\hline
				Overall wealth accumulation & Specific funds for the transaction \\
				\hline
				Inheritance, business profits, investments & Sale of property, business sale, salary \\
				\hline
				Broad picture & Specific transaction \\
				\hline
				Required for EDD (PEPs, high-risk) & Required for EDD (PEPs, high-risk) \\
				\hline
			\end{tabular}
		\end{block}
		\vspace{0.02cm}
		\begin{block}{\fontsize{7}{9}\selectfont \textbf{The TRAP}}
			\scriptsize
			\begin{itemize}
				\item \wronganswer{SOW and SOF are the same} — \textbf{NO.} Different concepts.
				\item \wronganswer{SOF is more important} — \textbf{BOTH} are required for EDD.
				\item \textbf{SOW} = overall wealth picture (how did you accumulate wealth?).
				\item \textbf{SOF} = where the specific funds came from (for this specific transaction).
				\item \textbf{BOTH} must be verified for high-risk clients.
				\item \textbf{Example:} SOW = inheritance from parents. SOF = sale of inherited property.
			\end{itemize}
		\end{block}
		\vspace{0.02cm}
		\recallbox{\scriptsize \textbf{Key Rule:} SOW = \textbf{overall wealth}. SOF = \textbf{specific funds}. Both required for EDD.}
	\end{frame}
	
	% ================================================================
	% SLIDE 12: OFAC - THE SANCTIONS TRAP
	% ================================================================
	\begin{frame}
		\frametitle{\fontsize{11}{13}\selectfont \textbf{OFAC — The Sanctions Trap}}
		\begin{block}{\fontsize{7}{9}\selectfont \textbf{The Hard Question — CAMS Exam Favorite}}
			\scriptsize
			What is the primary focus of the Office of Foreign Assets Control (OFAC) in U.S. financial regulation?
		\end{block}
		\vspace{0.02cm}
		\answerbox{\large \textbf{Issuing and enforcing economic and trade sanctions.}}
		\vspace{0.02cm}
		\begin{block}{\fontsize{7}{9}\selectfont \textbf{Why This Is a TRAP}}
			\scriptsize
			\begin{itemize}
				\item \wronganswer{Regulating anti-money laundering programs} — That's \textbf{FinCEN}, not OFAC.
				\item \wronganswer{Supervising banks} — That's \textbf{Federal Reserve/OCC}, not OFAC.
				\item \textbf{OFAC = sanctions enforcement.}
				\item \textbf{SDN List:} Specially Designated Nationals and Blocked Persons.
				\item \textbf{50\% Rule:} If sanctioned entity owns 50\%+ of another entity, that entity is also blocked.
				\item \textbf{Strict liability:} No intent required — if you process a transaction with a sanctioned entity, you're liable.
				\item \textbf{Key distinction:} OFAC = sanctions. FinCEN = AML/BSA.
			\end{itemize}
		\end{block}
		\vspace{0.02cm}
		\recallbox{\scriptsize \textbf{Key Rule:} OFAC = \textbf{sanctions enforcement}. FinCEN = \textbf{AML/BSA}. Don't confuse them.}
	\end{frame}
	
	% ================================================================
	% SLIDE 13: PSP - THE CUSTOMIZATION TRAP
	% ================================================================
	\begin{frame}
		\frametitle{\fontsize{11}{13}\selectfont \textbf{PSP — The Customization Trap}}
		\begin{block}{\fontsize{7}{9}\selectfont \textbf{The Scenario — CAMS Exam Favorite}}
			\scriptsize
			A PSP is planning to offer new services targeting small merchants in multiple countries, some of which lack mature AML frameworks.
		\end{block}
		\vspace{0.02cm}
		\begin{block}{\fontsize{7}{9}\selectfont \textbf{The Hard Question}}
			\scriptsize
			Which strategy is most appropriate to maintain compliance and mitigate risk?
		\end{block}
		\vspace{0.02cm}
		\answerbox{\large \textbf{Customize AML policy and onboarding procedures per jurisdiction and product.}}
		\vspace{0.02cm}
		\begin{block}{\fontsize{7}{9}\selectfont \textbf{Why This Is a TRAP}}
			\scriptsize
			\begin{itemize}
				\item \wronganswer{Apply home-country AML standards across all jurisdictions} — \textbf{NO.} One-size-fits-all fails.
				\item \textbf{Jurisdiction-specific:} High-risk jurisdictions = EDD, transaction limits, or restricted services.
				\item \textbf{Product-specific:} Low-value vs high-value products have different risk profiles.
				\item \textbf{Risk-based approach:} Apply proportionate measures based on risk.
				\item \textbf{FATF Rec 1:} Risk-based approach requires tailoring to specific risk profiles.
			\end{itemize}
		\end{block}
		\vspace{0.02cm}
		\recallbox{\scriptsize \textbf{Key Rule:} PSPs must \textbf{customize} AML per jurisdiction and product. One-size-fits-all fails.}
	\end{frame}
	
	% ================================================================
	% SLIDE 14: HIGH-VALUE ASSETS - THE CRYPTO RED FLAG TRAP
	% ================================================================
	\begin{frame}
		\frametitle{\fontsize{11}{13}\selectfont \textbf{High-Value Assets — The Crypto Red Flag Trap}}
		\begin{block}{\fontsize{7}{9}\selectfont \textbf{The Hard Question — CAMS Exam Favorite}}
			\scriptsize
			Which red flag is most commonly associated with abuse of high-value assets for laundering purposes?
		\end{block}
		\vspace{0.02cm}
		\answerbox{\large \textbf{Making payments for high-value items in cryptocurrency.}}
		\vspace{0.02cm}
		\begin{block}{\fontsize{7}{9}\selectfont \textbf{Why This Is a TRAP}}
			\scriptsize
			\begin{itemize}
				\item \wronganswer{Purchasing assets with complex financing plans spanning several years} — This \textbf{can} be a red flag, but crypto is \textbf{more directly} associated with laundering.
				\item \textbf{Cryptocurrency offers pseudonymity}, making it attractive for illicit payments.
				\item \textbf{High-value items:} Real estate, art, luxury cars purchased with crypto = suspicion.
				\item \textbf{Challenge:} Volatility and price fluctuations complicate source-of-funds verification.
				\item \textbf{Additional red flags:} Foreign PEPs purchasing multiple properties via trust companies.
				\item \textbf{Complex financing} = could be legitimate; \textbf{crypto payments} = more direct ML indicator.
			\end{itemize}
		\end{block}
		\vspace{0.02cm}
		\recallbox{\scriptsize \textbf{Key Rule:} Crypto payments for high-value assets = \textbf{primary red flag}.}
	\end{frame}
	
	% ================================================================
	% SLIDE 15: CASINO AML CONTROLS - THE SELECTION TRAP
	% ================================================================
	\begin{frame}
		\frametitle{\fontsize{11}{13}\selectfont \textbf{Casino AML Controls — The Selection Trap}}
		\begin{block}{\fontsize{7}{9}\selectfont \textbf{The Hard Question — CAMS Exam Favorite}}
			\scriptsize
			Which control measures help mitigate AML risks in casinos? (Select Two.)
		\end{block}
		\vspace{0.02cm}
		\answerbox{\large \textbf{1. Enforcing chip redemption policies tied to minimum gaming thresholds. 2. Tracking player activity through registered loyalty programs.}}
		\vspace{0.02cm}
		\begin{block}{\fontsize{7}{9}\selectfont \textbf{Why This Is a TRAP}}
			\scriptsize
			\begin{itemize}
				\item \wronganswer{Monitoring unstructured table gameplay for win/loss patterns} — This is a \textbf{monitoring technique}, not a formal control measure.
				\item \textbf{Chip redemption policies:} Prevents "chip washing" — exchanging chips for cash without meaningful gaming activity.
				\item \textbf{Thresholds:} US \$10,000 triggers reporting.
				\item \textbf{Loyalty programs:} Provide rich data on player behavior (wins, losses, session durations).
				\item \textbf{Enables:} Tracking of anomalies and suspicious patterns.
				\item \textbf{Unstructured monitoring} = too granular and not a formal control.
			\end{itemize}
		\end{block}
		\vspace{0.02cm}
		\recallbox{\scriptsize \textbf{Key Rule:} Casino controls = \textbf{chip redemption policies} + \textbf{loyalty program tracking}.}
	\end{frame}
	
	% ================================================================
	% SLIDE 16: E-COMMERCE MONITORING - THE VOLUME TRAP
	% ================================================================
	\begin{frame}
		\frametitle{\fontsize{11}{13}\selectfont \textbf{E-Commerce Monitoring — The Volume Trap}}
		\begin{block}{\fontsize{7}{9}\selectfont \textbf{The Hard Question — CAMS Exam Favorite}}
			\scriptsize
			Why is transaction monitoring particularly challenging in an e-commerce risk context?
		\end{block}
		\vspace{0.02cm}
		\answerbox{\large \textbf{High transaction volumes and decentralized payments make patterns harder to detect.}}
		\vspace{0.02cm}
		\begin{block}{\fontsize{7}{9}\selectfont \textbf{Why This Is a TRAP}}
			\scriptsize
			\begin{itemize}
				\item \wronganswer{Transactions are always small and therefore low-risk} — \textbf{NO.} Small transactions can still be suspicious.
				\item \textbf{High volume:} Millions of transactions daily — traditional rules-based monitoring generates too many false positives.
				\item \textbf{Decentralized payments:} Multiple payment gateways, cryptocurrencies, digital wallets.
				\item \textbf{Consequence:} Patterns are harder to detect.
				\item \textbf{Mitigation:}
				\begin{itemize}
					\item Use AI/ML models to detect anomalies
					\item Segment data by merchant, payment method, product
					\item Implement network analysis to identify patterns
				\end{itemize}
				\item \textbf{Additional risks:} Goods/services as ML vehicles (high-value goods purchased with illicit funds).
			\end{itemize}
		\end{block}
		\vspace{0.02cm}
		\recallbox{\scriptsize \textbf{Key Rule:} E-commerce challenge = \textbf{high volume + decentralized payments} = patterns harder to detect.}
	\end{frame}
	
	% ================================================================
	% SLIDE 17: DATA LINEAGE - THE TRANSPARENCY TRAP
	% ================================================================
	\begin{frame}
		\frametitle{\fontsize{11}{13}\selectfont \textbf{Data Lineage — The Transparency Trap}}
		\begin{block}{\fontsize{7}{9}\selectfont \textbf{The Hard Question — CAMS Exam Favorite}}
			\scriptsize
			Which support transparency and auditability through data lineage? (Select Two.)
		\end{block}
		\vspace{0.02cm}
		\answerbox{\large \textbf{1. Maintaining history of data transformations. 2. Capturing source of each data field.}}
		\vspace{0.02cm}
		\begin{block}{\fontsize{7}{9}\selectfont \textbf{Why This Is a TRAP}}
			\scriptsize
			\begin{itemize}
				\item \wronganswer{Archiving analyst notes from alert reviews} — This is \textbf{case documentation}, NOT data lineage.
				\item \wronganswer{Deleting outdated data annually} — \textbf{Undermines} auditability.
				\item \textbf{Data lineage:} Tracking the \textbf{flow of data} from source to destination.
				\item \textbf{History of transformations:} Shows how data was modified.
				\item \textbf{Source of each field:} Shows where data originated.
				\item \textbf{Why it matters:}
				\begin{itemize}
					\item Regulatory audit: Regulators want to see where data came from
					\item Model validation: Need to trace data inputs to models
					\item Error correction: Trace errors back to source
					\item SAR filing: Need to verify accuracy of reported data
				\end{itemize}
			\end{itemize}
		\end{block}
		\vspace{0.02cm}
		\recallbox{\scriptsize \textbf{Key Rule:} Data lineage = \textbf{transformations history} + \textbf{source of each field}. NOT case notes.}
	\end{frame}
	
	% ================================================================
	% SLIDE 18: DATA VALIDATION - THE OBJECTIVES TRAP
	% ================================================================
	\begin{frame}
		\frametitle{\fontsize{11}{13}\selectfont \textbf{Data Validation — The Objectives Trap}}
		\begin{block}{\fontsize{7}{9}\selectfont \textbf{The Hard Question — CAMS Exam Favorite}}
			\scriptsize
			Which are typical objectives of data validation and testing in AFC systems? (Select Three.)
		\end{block}
		\vspace{0.02cm}
		\answerbox{\large \textbf{1. Ensuring data inputs reflect actual customer and transaction behaviors. 2. Verifying that automated controls are functioning correctly based on input quality. 3. Testing whether controls generate alerts under appropriate risk conditions.}}
		\vspace{0.02cm}
		\begin{block}{\fontsize{7}{9}\selectfont \textbf{Why This Is a TRAP}}
			\scriptsize
			\begin{itemize}
				\item \wronganswer{Confirming reports are generated in a format preferred by third-party vendors} — \textbf{NOT} a validation objective.
				\item \textbf{Data inputs reflect actual behavior:} Test that data from source systems accurately feeds into AML systems.
				\item \textbf{Controls functioning correctly:} Validate that transaction monitoring rules fire alerts correctly.
				\item \textbf{Alerts under risk conditions:} Use known test cases to validate system logic.
				\item \textbf{Key principle:} Data validation ensures \textbf{accuracy} and \textbf{effectiveness} of AML controls.
				\item \textbf{Vendor preference is NOT} a regulatory or control objective.
			\end{itemize}
		\end{block}
		\vspace{0.02cm}
		\recallbox{\scriptsize \textbf{Key Rule:} Data validation = \textbf{accuracy of inputs} + \textbf{control functionality} + \textbf{appropriate alerts}.}
	\end{frame}
	
	% ================================================================
	% SLIDE 19: CLUSTERING - THE BEHAVIORAL TRAP
	% ================================================================
	\begin{frame}
		\frametitle{\fontsize{11}{13}\selectfont \textbf{Clustering — The Behavioral Trap}}
		\begin{block}{\fontsize{7}{9}\selectfont \textbf{The Hard Question — CAMS Exam Favorite}}
			\scriptsize
			Which are characteristics of clustering as used in compliance analytics? (Select Two.)
		\end{block}
		\vspace{0.02cm}
		\answerbox{\large \textbf{1. It groups entities by shared behavioral characteristics. 2. It identifies connections across multiple transaction types.}}
		\vspace{0.02cm}
		\begin{block}{\fontsize{7}{9}\selectfont \textbf{Why This Is a TRAP}}
			\scriptsize
			\begin{itemize}
				\item \wronganswer{It excludes outliers to reduce financial crime risk} — \textbf{NO.} Outliers are often the \textbf{most suspicious} — excluding them would be counterproductive.
				\item \textbf{Clustering:} Groups entities that exhibit similar behaviors.
				\item \textbf{Purpose:} Detect networks of criminal activity (money mule networks, organized crime).
				\item \textbf{Connections across transaction types:} Links separate accounts that exhibit similar patterns.
				\item \textbf{Does NOT:} Eliminate false positives, file SARs automatically, or replace manual review.
				\item \textbf{Does:} Reveal shared patterns among entities.
			\end{itemize}
		\end{block}
		\vspace{0.02cm}
		\recallbox{\scriptsize \textbf{Key Rule:} Clustering = \textbf{shared behavioral characteristics} + \textbf{connections across transaction types}.}
	\end{frame}
	
	% ================================================================
	% SLIDE 20: DIGITAL ONBOARDING - THE BENEFITS TRAP
	% ================================================================
	\begin{frame}
		\frametitle{\fontsize{11}{13}\selectfont \textbf{Digital Onboarding — The Benefits Trap}}
		\begin{block}{\fontsize{7}{9}\selectfont \textbf{The Hard Question — CAMS Exam Favorite}}
			\scriptsize
			Which benefits are associated with digital onboarding technology according to FATF? (Select Two.)
		\end{block}
		\vspace{0.02cm}
		\answerbox{\large \textbf{1. Lower cost and increased auditability. 2. Greater accountability in compliance processes.}}
		\vspace{0.02cm}
		\begin{block}{\fontsize{7}{9}\selectfont \textbf{Why This Is a TRAP}}
			\scriptsize
			\begin{itemize}
				\item \wronganswer{Elimination of human review requirements for high-risk customers} — \textbf{NO.} High-risk customers \textbf{still require} human review and EDD.
				\item \textbf{Lower cost:} Automation reduces manual effort.
				\item \textbf{Increased auditability:} Digital records are easier to track and audit.
				\item \textbf{Greater accountability:} Clear audit trails of who did what.
				\item \textbf{NOT:} Elimination of human review — human judgment is still required, especially for high-risk cases.
				\item \textbf{FATF:} Recognizes digital onboarding as beneficial for efficiency and inclusion.
			\end{itemize}
		\end{block}
		\vspace{0.02cm}
		\recallbox{\scriptsize \textbf{Key Rule:} Digital onboarding = \textbf{lower cost} + \textbf{auditability} + \textbf{accountability}. NOT elimination of human review.}
	\end{frame}
	
	% ================================================================
	% SLIDE 21: PERPETUAL KYC - THE TRIGGER-BASED TRAP
	% ================================================================
	\begin{frame}
		\frametitle{\fontsize{11}{13}\selectfont \textbf{Perpetual KYC — The Trigger-Based Trap}}
		\begin{block}{\fontsize{7}{9}\selectfont \textbf{The Hard Question — CAMS Exam Favorite}}
			\scriptsize
			Which feature most effectively distinguishes perpetual KYC from traditional periodic KYC reviews?
		\end{block}
		\vspace{0.02cm}
		\answerbox{\large \textbf{Trigger-based updates rather than fixed calendar reviews.}}
		\vspace{0.02cm}
		\begin{block}{\fontsize{7}{9}\selectfont \textbf{Why This Is a TRAP}}
			\scriptsize
			\begin{itemize}
				\item \wronganswer{It eliminates all manual review} — \textbf{NO.} Human review is still required.
				\item \wronganswer{It requires less customer data} — \textbf{NO.} It requires the same data, updated more frequently.
				\item \textbf{Perpetual KYC (pKYC):} Continuous, event-driven KYC refresh.
				\item \textbf{Triggers:}
				\begin{itemize}
					\item Changes in PEP status
					\item Adverse media hits
					\item Unusual transaction patterns
					\item Jurisdictional risk changes
					\item Significant changes in customer profile
				\end{itemize}
				\item \textbf{Benefits:}
				\begin{itemize}
					\item Reduces customer friction (fewer repeated KYC requests)
					\item Improves accuracy (updates in real-time)
					\item Ensures timely EDD (when risk changes)
				\end{itemize}
			\end{itemize}
		\end{block}
		\vspace{0.02cm}
		\recallbox{\scriptsize \textbf{Key Rule:} pKYC = \textbf{trigger-based updates}, not fixed calendar reviews.}
	\end{frame}
	
	% ================================================================
	% SLIDE 22: DEEPFAKE DETECTION - THE LIVENESS TRAP
	% ================================================================
	\begin{frame}
		\frametitle{\fontsize{11}{13}\selectfont \textbf{Deepfake Detection — The Liveness Trap}}
		\begin{block}{\fontsize{7}{9}\selectfont \textbf{The Scenario — CAMS Exam Favorite}}
			\scriptsize
			A fraudster used a high-resolution image and a deepfake video to bypass facial recognition verification.
		\end{block}
		\vspace{0.02cm}
		\begin{block}{\fontsize{7}{9}\selectfont \textbf{The Hard Question}}
			\scriptsize
			What tools could have prevented this fraud attempt? (Select Two.)
		\end{block}
		\vspace{0.02cm}
		\answerbox{\large \textbf{1. Use of software that can detect inconsistencies in light or skin texture. 2. Prompts for the customer to blink, smile, or repeat a phrase.}}
		\vspace{0.02cm}
		\begin{block}{\fontsize{7}{9}\selectfont \textbf{Why This Is a TRAP}}
			\scriptsize
			\begin{itemize}
				\item \wronganswer{Real-time fraud risk scoring algorithms using behavioral analytics} — This is \textbf{behavioral analytics}, not deepfake detection.
				\item \textbf{Texture analysis:} Deepfakes often have inconsistencies in light, skin texture, and shadows.
				\item \textbf{Active liveness prompts:} Blink, smile, repeat a phrase — requires real-time action.
				\item \textbf{Passive liveness:} Analyzes subtle facial movements without user action.
				\item \textbf{Active liveness is MORE secure} against deepfakes than passive.
				\item \textbf{Key distinction:} Behavioral analytics = monitoring behavior; liveness detection = verifying live presence.
			\end{itemize}
		\end{block}
		\vspace{0.02cm}
		\recallbox{\scriptsize \textbf{Key Rule:} Deepfake detection = \textbf{texture analysis} + \textbf{active liveness prompts}.}
	\end{frame}
	
	% ================================================================
	% SLIDE 23: AML COMMUNICATIONS - THE SCRIPT TRAP
	% ================================================================
	\begin{frame}
		\frametitle{\fontsize{11}{13}\selectfont \textbf{AML Communications — The Script Trap}}
		\begin{block}{\fontsize{7}{9}\selectfont \textbf{The Hard Question — CAMS Exam Favorite}}
			\scriptsize
			What should AML-sensitive communications materials emphasize to reduce risk? (Select Two.)
		\end{block}
		\vspace{0.02cm}
		\answerbox{\large \textbf{1. Avoid mentioning any internal reviews, even if prompted by the customer. 2. Use pre-approved scripts when addressing service disruptions caused by internal reviews.}}
		\vspace{0.02cm}
		\begin{block}{\fontsize{7}{9}\selectfont \textbf{Why This Is a TRAP}}
			\scriptsize
			\begin{itemize}
				\item \wronganswer{Refer customers to the AML department for financial advice} — \textbf{NO.} AML department doesn't give financial advice.
				\item \textbf{Avoid mentioning internal reviews:} This prevents tipping off.
				\item \textbf{Pre-approved scripts:} Ensures consistent, compliant messaging.
				\item \textbf{Tipping off risk:} Even hinting at an investigation can be tipping off.
				\item \textbf{Key principle:} Customer-facing staff must \textbf{not} reveal any AML-related investigations.
				\item \textbf{Best practice:} Have scripts for common scenarios (account holds, transaction delays).
			\end{itemize}
		\end{block}
		\vspace{0.02cm}
		\recallbox{\scriptsize \textbf{Key Rule:} AML communications = \textbf{avoid mentioning reviews} + \textbf{use pre-approved scripts}.}
	\end{frame}
	
	% ================================================================
	% SLIDE 24: ENTITY RESOLUTION - THE BLOCKCHAIN TRAP
	% ================================================================
	\begin{frame}
		\frametitle{\fontsize{11}{13}\selectfont \textbf{Entity Resolution — The Blockchain Trap}}
		\begin{block}{\fontsize{7}{9}\selectfont \textbf{The Hard Question — CAMS Exam Favorite}}
			\scriptsize
			How does entity resolution improve blockchain analytics for on-chain behavior monitoring?
		\end{block}
		\vspace{0.02cm}
		\answerbox{\large \textbf{It links data from multiple sources to improve match accuracy.}}
		\vspace{0.02cm}
		\begin{block}{\fontsize{7}{9}\selectfont \textbf{Why This Is a TRAP}}
			\scriptsize
			\begin{itemize}
				\item \wronganswer{It replaces all manual reviews} — \textbf{NO.} It's a tool, not a replacement.
				\item \wronganswer{It eliminates false positives entirely} — \textbf{NO.} It reduces them but doesn't eliminate them.
				\item \textbf{Entity resolution:} Links on-chain wallet addresses to off-chain risk data.
				\item \textbf{Multiple sources:} Wallet addresses, KYC data, sanctions lists, PEP lists, adverse media.
				\item \textbf{Example:} A wallet address associated with ransomware is linked to an exchange account.
				\item \textbf{NOT just name matching:} It's about linking \textbf{multiple identifiers} across systems.
				\item \textbf{Improves match accuracy:} Reduces false positives and false negatives.
			\end{itemize}
		\end{block}
		\vspace{0.02cm}
		\recallbox{\scriptsize \textbf{Key Rule:} Entity resolution = \textbf{links data from multiple sources} to improve match accuracy.}
	\end{frame}
	
	% ================================================================
	% SLIDE 25: SAR/STR - THE THRESHOLD TRAP
	% ================================================================
	\begin{frame}
		\frametitle{\fontsize{11}{13}\selectfont \textbf{SAR/STR — The Threshold Trap}}
		\begin{block}{\fontsize{7}{9}\selectfont \textbf{The Hard Question — CAMS Exam Favorite}}
			\scriptsize
			Is there a minimum amount for SAR/STR filing?
		\end{block}
		\vspace{0.02cm}
		\answerbox{\large \textbf{No minimum amount — file based on reasonable suspicion.}}
		\vspace{0.02cm}
		\begin{block}{\fontsize{7}{9}\selectfont \textbf{Why This Is a TRAP}}
			\scriptsize
			\begin{itemize}
				\item \wronganswer{\$5,000 minimum} — That's a \textbf{US SAR threshold} (FIs), but NOT universal.
				\item \wronganswer{\$10,000 minimum} — That's \textbf{CTR}, not SAR.
				\item \textbf{SAR/STR = reasonable suspicion}, NOT amount.
				\item \textbf{Attempted transactions:} Must be reported (even if not completed).
				\item \textbf{Structuring:} Must be reported (even if small amounts).
				\item \textbf{UK STR:} No minimum.
				\item \textbf{US SAR:} \$5,000+ for most FIs, but suspicious activity below threshold may still be filed.
				\item \textbf{Key principle:} Suspicion drives filing, NOT amount.
			\end{itemize}
		\end{block}
		\vspace{0.02cm}
		\recallbox{\scriptsize \textbf{Key Rule:} SAR/STR = \textbf{reasonable suspicion}, NOT amount. No minimum required.}
	\end{frame}
	
	% ================================================================
	% SLIDE 26: FINAL EXAM TIPS - ALL DOMAINS
	% ================================================================
	\begin{frame}
		\frametitle{\fontsize{11}{13}\selectfont \textbf{ALL DOMAINS — Hardest Exam Traps Summary}}
		\scriptsize
		\begin{block}{\fontsize{7}{9}\selectfont \textbf{Top 20 Traps — MUST KNOW FOR CAMS}}
			\scriptsize
			\begin{enumerate}
				\item \textbf{VASP/MiCA:} White paper + one notification, NOT each country.
				\item \textbf{Authentication:} Password biometrics = TRAP. Use MFA + physiological + behavioral.
				\item \textbf{MLAT:} Challenge is \textbf{admissibility}, not just sharing.
				\item \textbf{PPP:} Real-time two-way communication, NOT post-hoc anonymized reporting.
				\item \textbf{EU Directives:} Failure = \textbf{poor transposition}, not lack of definitions.
				\item \textbf{AMLA:} Cross-border, high-risk = AMLA, not national regulators.
				\item \textbf{DORA:} Centralized ICT oversight, NOT just national cybersecurity.
				\item \textbf{TCSP:} Cannot outsource AML responsibility. Decline unless BO disclosed.
				\item \textbf{PEP EDD:} Foreign = mandatory. Domestic = risk-based.
				\item \textbf{SOW vs SOF:} SOW = overall wealth. SOF = specific funds. Both required.
				\item \textbf{OFAC:} Sanctions enforcement. FinCEN = AML/BSA.
				\item \textbf{PSP:} Customize per jurisdiction, NOT one-size-fits-all.
				\item \textbf{High-Value Assets:} Crypto payments = primary red flag.
				\item \textbf{Casino:} Controls = chip redemption + loyalty program tracking.
				\item \textbf{E-Commerce:} Challenge = high volume + decentralized payments.
				\item \textbf{Data Lineage:} Transformations history + source of each field.
				\item \textbf{Data Validation:} Accuracy + control functionality + appropriate alerts.
				\item \textbf{Clustering:} Shared behaviors + connections across transaction types.
				\item \textbf{Digital Onboarding:} Lower cost + auditability. NOT elimination of human review.
				\item \textbf{Perpetual KYC:} Trigger-based updates, NOT fixed calendar reviews.
			\end{enumerate}
		\end{block}
		\vspace{0.02cm}
		\recallbox{\scriptsize \textbf{Final Rule:} Understand the \textbf{why} behind each concept. The exam tests \textbf{application}, not memorization.}
	\end{frame}
	
	% ================================================================
	% END SLIDE
	% ================================================================
	\begin{frame}
		\centering
		\vspace{1.5cm}
		{\fontsize{18}{22}\selectfont \textbf{ALL CAMS DOMAINS — COMPLETE}}\\[0.2cm]
		{\fontsize{11}{13}\selectfont Hard Questions \& Exam Traps from All Images}\\[0.1cm]
		\rule{4cm}{0.5pt}\\[0.1cm]
		{\fontsize{8}{10}\selectfont Domain A: VASP · Crypto · FATF Recs · Sanctions · DORA}\\[0.05cm]
		{\fontsize{8}{10}\selectfont Domain B: MLAT · PPP · FATF IOs · AMLA · Three Lines}\\[0.05cm]
		{\fontsize{8}{10}\selectfont Domain C: AML Policy · Board Oversight · PEP EDD · SOW/SOF}\\[0.05cm]
		{\fontsize{8}{10}\selectfont Domain D: Monitoring · OSINT · Data Lineage · AI/ML · Clustering}\\[0.1cm]
		\rule{4cm}{0.5pt}\\[0.1cm]
		{\fontsize{8}{10}\selectfont \textit{All traps explained. All concepts covered. Ready for the exam.}}
		\vspace{0.5cm}
	\end{frame}
	
\end{document}